\documentclass[10pt,a4paper]{amsart}
\usepackage[margin=19mm]{geometry}
\usepackage{amsmath,amssymb,mathtools}
\usepackage[scr=rsfs,cal=euler]{mathalfa}
\usepackage{xfrac}
\usepackage[unicode,hidelinks,bookmarks=false]{hyperref}
\newcommand{\ie}{i.e.\ }
\newcommand{\cf}{cf.\ }
\newcommand{\resp}{respectively\ }
\newcommand{\Subsection}{\subsection}
\providecommand{\nobreakdash}{\nobreak\hbox{-}}
\setlength{\emergencystretch}{2em}
\multlinegap0pt
\usepackage{mathtools}
\usepackage{amssymb}
\usepackage[shortlabels]{enumitem}
\usepackage{csquotes}
\newtheorem{lemma}{Lemma}
\newtheorem{corollary}{Corollary}
\newcommand{\R}{\mathbb R}
\title{An improved upper bound for the Froude number of irrotational solitary water waves}
\author{Evgeniy Lokharu, J\"org Weber}
\date{Source: arXiv:2502.18181v2 (CC BY 4.0)}
\begin{document}
\maketitle

\noindent\emph{Source and licence.} An improved upper bound for the Froude number of irrotational solitary water waves. Version arXiv:2502.18181v2 (CC BY 4.0), distributed by arXiv under the Creative Commons Attribution 4.0 International licence, http://creativecommons.org/licenses/by/4.0/, as recorded in the arXiv licence metadata for this exact version and shown on the abstract page https://arxiv.org/abs/2502.18181v2. Full licence notice: \texttt{licenses/arxiv-2502.18181-CC-BY-4.0.txt}.\par\smallskip

\noindent\textit{Editorial scope.} Counted unit: the article's corollary cor:u\_crestline (source0.tex lines 256--270). The article's complete original proof of it is delivered with it (source0.tex lines 271--292, one \texttt{proof} environment of 22 lines). No mathematics is written, repaired or re-derived: the statement and the proof are the article's own lines, in the article's own notation, and no retained line is substituted or re-flowed.  Retained uncounted as the source-local context the proof needs: lines 205--211.  Nothing was omitted from any retained span, so the package carries no omission record.  No retained line contains a citation, so the wrapper carries no bibliography and no bibliography entry.  No retained line contains a figure environment, an image-inclusion command or drawing code, so no image is delivered and the package declares no asset dependency.  Editorial formatting only, in addition: an AMS \texttt{amsart} preamble that restates the article's own declarations and macros used by the retained lines, namely the environments \texttt{lemma}, \texttt{corollary} on separate counters, together with the standard packages those lines need.  The retained environments print in this document as Lemma 1, Corollary 1 in order of appearance, whereas the article prints the same items with the numbering of its own full document.  The complete native source of the article as distributed in the arXiv e-print for this exact version is bundled unchanged as \texttt{sources/173834/source0.tex}.
\begin{lemma} \label{thm:slope}
	Given an arbitrary solitary wave solution, consider the harmonic function 
	\[
	f \coloneqq (\psi_y^2-\psi_x^2)\psi_{xy}-2\psi_x\psi_y\psi_{xx}+\tfrac35\psi_x = -(u^2-v^2)u_x+2u u_y v+\tfrac35v.
	\]
	Then, $f$ is negative everywhere in $\Omega\coloneqq\{(x,y)\in\R^2:0<y<\eta(x),-\infty<x<0\}$, which is the left open half of the fluid domain, where $\eta' > 0$.
\end{lemma}

\begin{corollary}\label{cor:u_crestline}
	Everywhere below the crest we have
	\begin{equation}\label{eq:cor_monotone}
		(\psi_y^2 \psi_{yy}+\tfrac35 \psi_y)_y = -(u^2 u_y+\tfrac35 u)_y < 0.
	\end{equation}
	Moreover, at the crest line,
	\begin{equation}\label{eq:est_psiy_up_good}
		-u = \psi_y \le \frac{\sqrt2\sqrt{80r-27y-53\hat\eta}}{3\sqrt5} - \frac{\sqrt{2(r-\hat\eta)}}{3}.
	\end{equation}
	In particular,
	\begin{equation}\label{eq:est_psiy_up_good_extreme}
		u^2 \leq \tfrac65 (r-y)
	\end{equation}
	at the crest line for an extreme wave.
\end{corollary}

\begin{proof}
	First, \eqref{eq:cor_monotone} follows directly from the Hopf lemma, since $f = 0$ below the crest, while $0$ is the global maximum of $f$ on the closure of $\Omega$.
	
	Let us now consider an extreme wave for which $\psi_y(0,\hat\eta)=0$. Then, by \eqref{eq:cor_monotone}, $\psi_y^2 \psi_{yy}+\tfrac35 \psi_y\ge0$ at the crest line. Dividing by $\psi_y$ and integrating from $y$ to $\hat\eta$ yields \eqref{eq:est_psiy_up_good_extreme}.
	
	Consider now a non-extreme wave. At the crest, we have $\psi_{yy}=\eta''\psi_y$ in view of $T^2\psi=0$. Similarly as in the previous proof, $\eta''$ can be expressed in terms of $\psi_y$, $\eta'$, and $f$:
	\[\eta'' = -\frac{5(1+\eta'^2)^2 f + 8\eta' (1-3\eta'^2+\eta'^4)\psi_y}{20\eta' (1-\eta'^2) \psi_y^3}.\]
	Since $f/\eta'\le0$ for $x\ne0$ by Lemma \ref{thm:slope}, we get
	\[\eta'' \ge -\frac{2}{5\psi_y^2}\]
	at the crest. Therefore, at the crest line,
	\[\psi_y^2 \psi_{yy}+\tfrac35 \psi_y \ge \frac{\psi_y(0,\hat\eta)}{5} > 0.\]
	Thus, we have a differential inequality of the form
	\[h' + ah^{1/3} \ge b,\]
	with $h=\psi_y^3(0,\cdot)$, $a=9/5$, and $b=3\psi_y(0,\hat\eta)/5=3\sqrt{2(r-\hat\eta)}/5$, such that $ah^{1/3}>b$. This inequality can be rearranged and integrated from $y$ to $\hat\eta$ to yield
	\[y-\hat\eta \le \left[\frac{3(2b+ah^{1/3})h^{1/3}}{2a^2} + \frac{3b^2\ln(ah^{1/3}-b)}{a^3}\right]\Bigg|_{h=h(y)}^{h(\hat\eta)}.\]
	Therefore,
	\begin{align*}
		\psi_y^2 + \tfrac23 \sqrt{2(r-\hat\eta)}\psi_y &\le \tfrac65 (\hat\eta-y) + \tfrac29 (r-\hat\eta) \left(15 + \ln\frac{8(r-\hat\eta)}{\left(3\psi_y-\sqrt{2(r-\hat\eta)}\right)^2}\right)\\
		&\le \tfrac65 (\hat\eta-y) + \tfrac{10}{3} (r-\hat\eta),
	\end{align*}
	where in the last step we simply used $\psi_y\ge\sqrt{2(r-\hat\eta)}$. Completing the square, we arrive at \eqref{eq:est_psiy_up_good}.
\end{proof}

\end{document}
